\documentclass[12pt]{article}
\usepackage[utf8]{inputenc}
\usepackage[T1]{fontenc}
\usepackage{hyperref}
\title{Yatzachi Zapotec Language Resources for K-12 Curriculum Development}
\author{Compiled by Claude Sonnet 5.5 (claude-sonnet-5-5) for Andr\'as Kornai}
\date{}
\begin{document}
\maketitle

\section{Introduction}
Yatzachi Zapotec (ISO 639-3 zav, Glottocode yatz1235), also called Zapotec, Villa Alta, is an Otomanguean language of Oaxaca, Mexico, with about 2,500 speakers according to a 1990 census figure cited by the Digital Bible Society. The resources found are mostly missionary and linguistic work from SIL, and nothing aimed at schools turned up.

\section{Primary Educational Resources}
No primers, textbooks or school programmes were found in the search results.

\section{Linguistic and Reference Materials}
Inez Butler's \emph{Diccionario Zapoteco de Yatzachi} is cited in several databases. The dates differ: 1997 in ASJP and 2000 in the Language Description Heritage catalogue. WALS cites Butler's 1980 \emph{Gram\'atica zapoteca} for the Yatzachi el Bajo variety. No online copy was found. The links below are catalogue or database entries.
\begin{itemize}
\item Language Description Heritage entry: \url{https://ldh.clld.org/languages/zav}
\item ASJP entry: \url{https://asjp.clld.org/languages/ZAPOTEC_YATZACHI}
\item WALS entry: \url{https://wals.info/languoid/lect/wals_code_zya}
\end{itemize}

\section{Religious and Community Resources}
A Yatzachi Zapotec New Testament exists. Sources disagree on its date: 1970 at the Digital Bible Society and 1971 at eBible.org.
\begin{itemize}
\item Digital Bible Society, Zapoteco de Yatzachi New Testament: \url{https://dbs.org/bibles/ZAVTBL}
\item eBible.org, details page: \url{https://ebible.org/find/details.php?id=zavtbl}
\item eBible.org, New Testament: \url{https://ebible.org/Scriptures/details.php?id=zavNT}
\item Scripture Earth, with text and audio: \url{https://www.scriptureearth.org/00eng.php?iso=zav}
\end{itemize}

\section{Digital and Audio Resources}
Global Recordings Network offers Bible-story audio and video, in MP3 and MP4 slideshow packages.
\begin{itemize}
\item Global Recordings Network, Yatzachi Zapotec: \url{https://globalrecordings.net/en/language/67}
\item Global Recordings Network, Zapotec, Villa Alta: \url{https://globalrecordings.net/en/language/zav}
\end{itemize}

\section{Curriculum Development Context}
The language is endangered and is not known to be taught in schools. Existing material is limited to religious texts and audio, plus older SIL linguistic descriptions. Butler's dictionary and grammar, and the narrative corpus said to exist, are the natural starting points. Whether they are accessible and how they are licensed would need to be checked.

\end{document}